\documentclass{cukthesis} % add [german] if required

\def\Title{Title} % Hier Titel eintragen / Enter your title here
\def\Author{Niklas Isermann} % Hier Namen eintrage / Enter your name here
\def\Degree{Bachelor} % Master
\def\SecondExaminer{???} % Zweitgutachter / Second Examiner

%Some useful default libs
\usepackage{graphicx}
\usepackage{url}
\usepackage{hyperref} %comment out if this causes trouble. Otherwise, this makes references in the document clickable.
\usepackage{todonotes} %allows for the \todo{...} and \todo[inline]{...} commands
\usepackage{cleveref} %use \Cref instead of \ref to reference labels (e.g. \label{def:ddh}). Allows for example \Cref{def:ddh} which is rendered as Definition X.Y. Documentation: https://mirror.physik.tu-berlin.de/pub/CTAN/macros/latex/contrib/cleveref/cleveref.pdf

% Add additional packages as required

%Some typical commands (remove or change if they're not to your liking)
\renewcommand{\epsilon}{\varepsilon}
\DeclareMathOperator*{\argmax}{arg\,max}
\DeclareMathOperator*{\argmin}{arg\,min}
\newcommand*{\bfrac}[2]{\genfrac{\lbrace}{\rbrace}{0pt}{}{#1}{#2}} %Fraction withour bar

\newcommand{\N}{\mathbb{N}}
\newcommand{\Z}{\mathbb{Z}}
\newcommand{\R}{\mathbb{R}}

%Example defs for the example.tex file
\newcommand{\setup}{\mathsf{Setup}}
\newcommand{\gobble}{\mathsf{Gobble}}
\newcommand{\ungobble}{\mathsf{Ungobble}}
\newcommand{\pk}{\mathit{pk}}

% Add additional style definitions as required

\definecolor{mygreen}{rgb}{0,0.6,0}
\definecolor{mygray}{rgb}{0.5,0.5,0.5}
\definecolor{mymauve}{rgb}{0.58,0,0.82}
\usepackage{amsmath}
\usepackage{hyperref}
\usepackage{comment}
\usepackage{listings}
\usepackage{pgfplots}
\usepackage{filecontents}
\usepackage{pgf-pie}
\pgfplotsset{compat = newest}
\usepackage{float}

\def\code#1{\texttt{#1}}
\lstset{ 
	backgroundcolor=\color{white},   % choose the background color; you must add \usepackage{color} or \usepackage{xcolor}; should come as last argument
	basicstyle=\footnotesize,        % the size of the fonts that are used for the code
	breakatwhitespace=false,         % sets if automatic breaks should only happen at whitespace
	breaklines=true,                 % sets automatic line breaking
	captionpos=b,                    % sets the caption-position to bottom
	commentstyle=\color{mygreen},    % comment style
	deletekeywords={...},            % if you want to delete keywords from the given language
	escapeinside={\%*}{*)},          % if you want to add LaTeX within your code
	extendedchars=true,              % lets you use non-ASCII characters; for 8-bits encodings only, does not work with UTF-8
	firstnumber=0,                % start line enumeration with line 1000
	frame=single,	                   % adds a frame around the code
	keepspaces=true,                 % keeps spaces in text, useful for keeping indentation of code (possibly needs columns=flexible)
	keywordstyle=\color{blue},       % keyword style
	language=Octave,                 % the language of the code
	morekeywords={*,...},            % if you want to add more keywords to the set
	numbers=left,                    % where to put the line-numbers; possible values are (none, left, right)
	numbersep=5pt,                   % how far the line-numbers are from the code
	numberstyle=\tiny\color{mygray}, % the style that is used for the line-numbers
	rulecolor=\color{black},         % if not set, the frame-color may be changed on line-breaks within not-black text (e.g. comments (green here))
	showspaces=false,                % show spaces everywhere adding particular underscores; it overrides 'showstringspaces'
	showstringspaces=false,          % underline spaces within strings only
	showtabs=false,                  % show tabs within strings adding particular underscores
	stepnumber=2,                    % the step between two line-numbers. If it's 1, each line will be numbered
	stringstyle=\color{mymauve},     % string literal style
	tabsize=2,	                   % sets default tabsize to 2 spaces
	title=\lstname                   % show the filename of files included with \lstinputlisting; also try caption instead of title
}
\begin{document}

\bibliographystyle{alpha}

\pagenumbering{roman}

\thispagestyle{empty}
\begin{titlepage}
\begin{center}

	\begin{minipage}{14cm}
		\hspace*{1.9cm}
		\includegraphics[height=3cm]{figures/upb_logo}\\
		\hspace*{1.15cm}
		\begin{minipage}{2.8cm}
			\includegraphics[width=2.5cm]{figures/cuk_logo_2019}
		\end{minipage}
		\begin{minipage}{9.8cm}
			\vspace*{5pt}
			\textsf{%
			Fakultät für Elektrotechnik, Informatik und Mathematik \\
			Arbeitsgruppe Codes und Kryptographie
			}
		\end{minipage}
	\end{minipage}\\[60pt]

	\begin{doublespace}
		{\Huge\textbf{\Title}}\\[30pt]
	\end{doublespace}

	{\Large
		\ifgerman
			\Degree arbeit
		\else
			\Degree 's Thesis
		\fi
	}\\[6pt]
		\ifgerman
			im Rahmen des Studiengangs Informatik\\
			zur Erlangung des Grades
		\else
			in Partial Fulfillment of the Requirements for the\\
			Degree of
		\fi
		\\[6pt]
  	{\Large \Degree\ of Science}\\[54pt] % Replace if different

	\ifgerman
		von\\
	\else
		by\\
	\fi
	{\scshape\large \Author}\\[54pt]

	\ifgerman
		vorgelegt bei:\\
	\else
		submitted to:\\
	\fi

	{\large Prof. Dr. Johannes Blömer \\
	\ifgerman
		und
	\else
		and
	\fi
	\\[6pt]
	\large \SecondExaminer}\\[30pt]

	{Paderborn, \today}

\end{center}
\end{titlepage}
\clearpage


\newgeometry{top=10mm, bottom=20mm, right=20mm}
\begin{small}

  \parindent0ex

  \chapter*{Eidesstattliche Versicherung}

  \thispagestyle{empty}

  \bigskip

  Nachname: \parbox[t]{.343\linewidth}{\rule[-3pt]{\linewidth}{.4pt}} \hspace{2ex} Vorname:  \parbox[t]{.35\linewidth}{\rule[-3pt]{\linewidth}{.4pt}}

  \vspace{2\bigskipamount}

  Matrikelnr.: \parbox[t]{.15\linewidth}{\rule[-3pt]{\linewidth}{.4pt}} \hspace{2ex} Studiengang: \parbox[t]{.5\linewidth}{\rule[-3pt]{\linewidth}{.4pt}}

  \bigskip

  $\square$ Bachelorarbeit \hspace{1ex} $\square$ Masterarbeit

  \bigskip

  Titel der Arbeit: Evaluating database systems relying on secure multi-party computation 

  \bigskip

  $\square$ Die elektronische Fassung ist der Abschlussarbeit beigefügt.

  \bigskip

  $\square$ Die elektronische Fassung sende ich an die/den erste/n Prüfenden bzw.~habe ich an die/den erste/n Prüfenden gesendet.

  \bigskip

  Ich versichere hiermit an Eides statt, dass ich die vorliegende Abschlussarbeit (Ausarbeitung inkl.~Tabellen, Zeichnungen, etc.)~selbstständig und ohne unzulässige fremde Hilfe erbracht habe.
  Ich habe keine anderen als die angegebenen Quellen und Hilfsmittel benutzt sowie wörtliche und sinngemäße Zitate kenntlich gemacht.
  Die Abschlussarbeit hat in gleicher oder ähnlicher Form noch keiner Prüfungsbehörde vorgelegen.
  Die elektronische Fassung entspricht der gedruckten und gebundenen Fassung.

  \subsection*{Belehrung}

  Wer vorsätzlich gegen eine die Täuschung über Prüfungsleistungen betreffende Regelung einer Hochschulprü\-fungsordnung verstößt, handelt ordnungswidrig.
  Die Ordnungswidrigkeit kann mit einer Geldbuße von bis zu 50.000,00 \euro~geahndet werden.
  Zuständige Verwaltungsbehörde für die Verfolgung und Ahndung von Ordnungswidrigkeiten ist die Vizepräsidentin / der Vizepräsident für Wirtschafts- und Personalverwaltung der Universität Paderborn.
  Im Falle eines mehrfachen oder sonstigen schwerwiegenden Täuschungsversuches kann der Prüfling zudem exmatrikuliert werden.~(§~63 Abs.~5 Hochschulgesetz NRW in der aktuellen Fassung).

  \medskip

  Die Universität Paderborn wird ggf.~eine elektronische Überprüfung der Abschlussarbeit durchführen, um eine Täuschung festzustellen.

  \medskip

  Ich habe die oben genannten Belehrungen gelesen und verstanden und bestätige dieses mit meiner Unterschrift.

  \vspace{2\bigskipamount}

  Ort: \parbox[t]{.4\linewidth}{\rule[-3pt]{\linewidth}{.4pt}} \hspace{2ex} Datum: \parbox[t]{.15\linewidth}{\rule[-3pt]{\linewidth}{.4pt}}

  \vspace{2\bigskipamount}

  Unterschrift: \parbox[t]{.5\linewidth}{\rule[-3pt]{\linewidth}{.4pt}}

  \subsection*{Datenschutzhinweis}

  Die o.g.~Daten werden aufgrund der geltenden Prüfungsordnung (Paragraph zur Abschlussarbeit) i.V.m.~§~63 Abs.~5 Hochschulgesetz NRW erhoben.
  Auf Grundlage der übermittelten Daten (Name, Vorname, Matrikelnummer, Studiengang, Art und Thema der Abschlussarbeit) wird bei Plagiaten bzw.~Täuschung der/die Prüfende und der Prüfungsausschuss Ihres Studienganges über Konsequenzen gemäß Prüfungsordnung i.V.m.~Hochschulgesetz NRW entscheiden.
  Die Daten werden nach Abschluss des Prüfungsverfahrens gelöscht.
  Eine Weiterleitung der Daten kann an die/den Prüfende/n und den Prüfungsausschuss erfolgen.
  Falls der Prüfungsausschuss entscheidet, eine Geldbuße zu verhängen, werden die Daten an die Vizepräsidentin für Wirtschafts- und Personalverwaltung weitergeleitet.
  Verantwortlich für die Verarbeitung im regulären Verfahren ist der Prüfungsausschuss Ihres Studiengangs der Universität Paderborn, für die Verfolgung und Ahndung der Geldbuße ist die Vizepräsidentin für Wirtschafts- und Personalverwaltung.

\end{small}

\restoregeometry


\tableofcontents

% If you have figures, uncomment this
%\listoffigures

% If you have algorithms, uncomment this
%\listofalgorithms

\cleardoublepage
\pagenumbering{arabic}

% Input your files here!
%\chapter{Basic definitions and notation}
\todo[inline]{This file contains example content. It is meant to get you started.}
\todo[inline]{You can remove this example content from the thesis by removing the input statement for example.tex from the thesis\_main.tex file.}

\section{Basic notation}
\label{sec:basicnotation}
Throughout this thesis, we will use the following notation:
\begin{itemize}
	\item $\N := \{1,2,\dots\}$ denotes the set of natural numbers (excluding zero). %see config/defs.tex file for the definition of the $\N$ macro.
	\item For a bit string $s = (s_0,\dots,s_{n-1}) \in \{0,1\}^n$ and $0\leq i \leq j < n$, we write $s[i:j] := (s_i, \dots, s_{j-1})$ to denote substrings.
	In particular, $s[i:i]$ is the empty string $\epsilon$ and $s[0:n]$ is the complete string $s$. 
	\item For two vectors $\vec{u},\vec{v}\in \{0,1\}^n$, with $\vec{u}=(u_1,\dots,u_n), \vec{v} = (v_1,\dots,v_n)$, the expression $\vec{u}\odot\vec{v}$ denotes the Hadamard product. $(\vec{u}\odot\vec{v})_i = u_i\cdot v_i$.
\end{itemize}

\section{Gobbling schemes}
\todo[inline]{Don't try to make sense of this. It's just a syntax example with no real connection to anything.}
Gobbling schemes \todo{This is an example for a todonote. They're super useful!} are useful whenever the the Hadamard product of two random bit vectors needs to be hidden from polynomially bounded adversaries. 
They are an important building block for twaddle signatures, which we investigate in this thesis.
\subsection{Syntax definition}
Our definition for gobbling schemes is taken from \cite{testref} \todo[fancyline]{note the citation} with minor syntactic changes. 
See Section~\ref{sec:basicnotation} \todo[fancyline]{note the reference} for basic notation.

\begin{definition}[Gobbling scheme]
	A \emph{gobbling scheme} $\Pi$ consists of the following three probabilistic polynomial-time algorithms:
	\begin{itemize}
		\item $\pk \leftarrow \setup(1^\lambda)$ \todo{note that $\pk$ and $\setup$ are macros defined in defs.tex.} 
			on input a unary security parameter $\lambda$, $\setup$ generates a public key $\pk$.
		\item $\vec{q} \leftarrow \gobble(\pk, \vec{u}, \vec{v})$ generates a \emph{gobbled} vector $\vec{q}\in\{0,1\}^n$ given a key $\pk$ and two vectors $\vec{u},\vec{v}\in\{0,1\}^n$.
		\item $\vec{z} \leftarrow \ungobble(\pk, \vec{q})$ given a gobbled vector $\vec{q}\in\{0,1\}^n$ outputs an ungobbled vector $\vec{z}\in\{0,1\}^{2n}$.
	\end{itemize}
	$\Pi$ is \emph{correct} if there exists a negligible function $\mu$ such that for all $\lambda,n\in\N$,
	\begin{align*}
		1-\mu(\lambda) \leq \Pr[\vec{z} = \vec{u}\odot\vec{v} \mid & k\leftarrow \setup(1^\lambda); \vec{u},\vec{v}\leftarrow \{0,1\}^n; \\
		& \vec{z} \leftarrow \ungobble(k, \gobble(k, \vec{u},\vec{v}))]
	\end{align*}
\end{definition}
\noindent Intuitively, correctness guarantees \dots

\subsection{Security definition}
\dots

\chapter {Introduction}
\section{MPC and Databases}
A famous problem in the context of MPC is Yao's millionaire's problem. In Yao's millionaire's problem there are two millionaires Alice and Bob. We will call Alice's wealth x and Bob's wealth y. Alice and Bob want to know who of them is has more money. i.e. they want to compute the function F(x,y) := $ 
\begin{dcases} 
	Alice \,  is \, richer   &  y \leq  x \\
	Bob \, is \, richer   &  y > x  \\
\end{dcases}  $. Yet neither of them is willing to trust the other and tell him how much money he has. 
Yao's millionaire's problem can be generalised into the general MPC problem.  Instead of Bob and Alice, we now consider n parties $ p_0,\dots,p_{n-1} $ and each party i holds an arbitrary input  $ x_i $ for an arbitrary function F($ x_0,\dots,x_{n-1}$),  that all parties have agreed upon.  
A MPC protocol $ \pi $  is protocol, that allows  $ p_0,\dots,p_{n-1} $ to compute F($ x_0,\dots,x_{n-1} )$ without revelling any information about $ x_0,\dots,x_{n-1}. $
       

             
Andrew Yao proposed a solution for Yao's millionaire's problem in 1982 [].  It has also been shown that MPC is Turing-complete[]. This means that for any function f that can be computed with a Turing machine. There exists a MPC protocol $ \pi $ that can compute f.    
-Databases ...










\chapter {related work}



\section{Contribution}
In our work we have evaluated existing secure multiparty computation frameworks that are able to perform typical database operations, like joins, union, and selects. We started with a pre-selection of 6 frameworks. First, we implemented different use cases with the frameworks on a trial basis, thereby we placed particular emphasis on the usability of the frameworks. Afterwards we selected the three frameworks that were most promising to yield good results and discarded the rest. We then defined several different use cases, that include typical database operations. We implemented the use cases with the selected frameworks and evaluated their performance. Thereby we used different technical tools, to keep our results as robust and reproducible as possible. Finally, we graphically present and discuss the collected data. 


\section{structure}
In this section we provide an overview over the rest of our work and summarise the content of each remaining chapter.
\paragraph{Preliminary}
In our second chapter we provide an overview over the topic of secure multiparty computation. We explain the general concept and different important security assumption. Furthermore we explain the database operations that are of importance four our work. 
\paragraph{Framework Description}
Our third chapter contains a detailed description of the frameworks we used to perform our benchmarks. Furthermore, a comparative discussion of the frameworks can be found in the chapter, where we especially focus on usability aspects. The chapter concludes with a description of the frameworks that did not make it into our final selection. 
\paragraph{Benchmarking}
In our fourth chapter, we describe the different metrics for which we collected data. In particular, we list the different tools that helped us to do this and provide an overview how they are used. 
\paragraph{Implementation}
In our fourth chapter, we describe our implantation. First, we describe our use cases and motivate their design.  Then we describe the concrete implementations with the individual frameworks, including commented excerpts from our source code.
\paragraph{Evaluation}
In our last chapter we evaluate the data we have collected. We highlight special features and provide a graphical representation of the data. 
Finally we conclude our work with a discussion of our results.

\chapter{Preliminary}
\section{Secure Multiparty Computation}
In the an secure multiparty computation(short MPC) scenario there are n parties $ p_0,\dots,p_{n-1} $. They want to compute an agreed upon functionality F($ x_0,\dots,x_{n-1} $). A functionality is a function that is allowed to have internal randomness, so its not function in the strict mathematically sense of the word. 
Each party $ p_i $ holds an input value $ x_i $. 
The parties hold their input private and do not want to reveal any information about it. The Goal of secure multiparty computation is to develop a protocol  $ \pi $ that enables them to jointly compute F($ x_0,\dots,x_{n-1}. $). The security goal of "not revelling the inputs" is often formalised through the Real-/Ideal-World Paradigm. 


\subsection{Real World and Ideal World}
When modelling security of secure multiparty computation we compare the real world, to a perfect ideal world, where the problem can be solved in a perfect way.
\paragraph{Real World}
In the real world there exists a protocol $\pi $ that enables the parties to compute F. All parties execute the protocol together. During the execution they exchange several rounds of communication. The attacker or adversary has the ability to corrupt one or more of the parties. His capability to influence the corrupted parties is an important parameter and may differ based on different security assumptions. These may range for example from a relative weak adversary that can only read massages to a very powerful adversary. We explain the adversary models that are of importance for our benchmarks in \hyperref[sec:Adversarial Models]{Section 3.2} in detail.
The real world view of party A consist of the input of A, all massages A sends or receives during the execution of the protocol and his internal randomness. The protocol achieves the security goal of confidentially if the attacker is unable to derive new information from the views of the parties he did corrupt. ...
\paragraph{Ideal World}\todo[fancyline]{Schaubild einfügen ?}
In the Ideal World the parties to not need run a protocol. Instead they can rely on a trusted, incorruptible third party P that aids them. With the aid of P, the parties can evaluate F in two simple steps. In a first round of communication every party send P its input. P now holds all information it needs to compute F. Afterwards P can send each party the result in a second round of communication. Like in the real world, in the ideal world their also exists an adversary. Similar to his real world counterparty he is also able to corrupt one or more parties.    Compared to his real world counter part the ideal world adversary has otherwise very limited ability's. He can only see the input and output of the parties he corrupts. Since the computation with the aid of P produces no intermediate results that he can observe. Depending on the underlying security assumptions he also may be able to modify the input a corrupted party sends to P in the first round of communication. Because of these very limited ability's it is desired that for every adversary in the real world their exits a comparable powerful adversary in the ideal world. This is often formalised using simulation based proof.

\paragraph{Security}
Given an real-world adversary A , a secure multiparty protocol $\pi $ , and a functionality F for $\pi $ to be secure we require the existence of a so called simulator S.
S is an ideal world adversary for F that mirror's the behaviour of A. This means that S and A corrupt the same parties and also that if a A changes the output of F then S does the same. After S has performed its attack S outputs a real world view. $\pi $ is secure against A, if the view S outputs is indistinguishable from a view of A. This means that the real world attacker A cannot learn more then the ideal world attacker S. Despite the very limited ability's S has compared to A. Finally we say $\pi $  is secure against if, for all A $\pi $ is secure against A. 






















     

\section{Adversarial Models}
\label{sec:Adversarial Models}
There are multiple models and categorizations of adversary's and their capability's. These distinctions have significant impact on feasibility and difficulty of secure multiparty computation. In the following we will outline the models and assumptions that are of importance for our benchmarks. 

\paragraph{Passive Adversary vs Active Adversary}
A passive adversary can not force a corrupted party to deviate from the protocol in an any way. One could think of a passive adversary as a "read-only" adversary. As a passive adversary is only able to read the messages his corrupted parties receive or send. A active adversary can do everything a passive adversary can additionally he has the power to force a corrupted party to deviate from the protocol in an arbitrary way. So if for example the protocol would at some point require that each party choses an integer between 1 and n uniformly at random. 
Then a passive adversary would have no choice but to choose the integer between 1 and n uniformly at random. 
On the contrary an active adversary would be able to force a corrupted party to chose the value 42 or any other value that the adversary considers to be advantageous for him. In the ideal world a passive adversary is bound to forward the real input values. A active adversary can choose to ignore the real input and forward any value instead.

\paragraph{Monolithic Adversary}
In the following we will assume a monolithic adversary unless explicitly stated otherwise. This means that there is only a single adversary that controls all corrupt parties. For the honest parties a monolithic adversary is a worst-case scenario. A monolithic adversary is more powerful compared to multiple adversary's that control the same total amount of parties but to not corporate with each other. A protocol that is secure in the presence of a single adversary that corrupts n parties and is able to coordinate their efforts. Will be secure in the presence of up to n adversary's that corrupt n parties total and do not coordinate their efforts.   
\paragraph{General Adversary vs Threshold Adversary}
In the threshold MPC setting the adversary can choose to corrupt any party. The threshold adversary is only limited in the way that can at most corrupt t parties where t is set to be 0 < t < n. A common setting for t is t = $\left \lfloor{ \frac{n}{2} }\right \rfloor  $, which called the honest majority. For example and for n=3 the presence of an honest majority means, that it is assumed that the threshold adversary can corrupt at most 1 party. Threshold MPC best fits scenarios that feature a very homogenous group of parties. A general adversary is limited in his choice which party he corrupts by an adversary structure  
$ Z = \{ Z_1, \dots, Z_l  \} $. Where $ Z_i $ can be any set of parties. The general adversary must corrupt a set of parties  $ P $ such that there exists an $ x \in Z $ that holds $ P \subset x $. This allows for a flexible way to formalise assumptions. If for example in protocol there are two parties that hold a very vital role and one want to assume that no adversary can corrupt both of these parties. That can be formalised by using a general adversary an defining $ Z $ so that no element in  $Z $ contains both of these two parties.  
 
\paragraph{Static vs Dynamic Corruptions}
Another important distinction is the distinction between static and adaptive adversary's. A static adversary is bound to chose which parties he wants to corrupt before the execution of $ \pi $ starts. An adaptive adversary can corrupt a party during the execution of the protocol. This makes the adaptive adversary much more powerful. As he can try to identify "weak links" based on the information he gets during the execution of the protocol and then choose corrupt those.   


\subsection{Additional Properties}
coming soon
\section{Databases}
coming soon
\chapter{framework description}

\section{conclave}
- published in 2019,\\
- combining MPC operations and cleartext processing\\ 
- moving operations outside of MPC to maximise performance\\
- allows to explicitly annotate trust between parties\\
- utilises trust through hybrid operations for additional performance increase \\
- maintaining same end-to-end security as "pure" MPC     \\
- relying on Obliv-C and Sharemind as MPC backend\\
- backend theoretic exchangeable through generic interface\\

- Query optimization aims to move as work outside of mpc\\
- contrary to conventional sql operations that aims to minimize the total amount of work e.g. filters before join \\

- secure against semi-honest adversary \\
- adversary is statically\\
- no secure channel setting\\
- inherits security guarantees of backend -> honest majority \\
- compares to "SMCQL most similar existing system"\\
- jiff dependency \\
- requires python 3.5 \\
- ...
\section{ABY3}
ABY3 ...
\subsection{functionality }
ABY3 is a 3-party MPC framework that allows to compute queries on relational database tables. It focuses on computing various SQL-like join as efficiently as possible. Therefore it features a large range of different join operations these include but are not limited to left join, right join, set union, set minus, and also full joins.  
Besides joins it is also possible to query a single table with query's that have a comparable semantic to the SELECT, FROM,WHERE; statement in SQL. For example a selection like "select X1 from X  where X2 > 42" can be done with relative ease using the implemented features of ABY3. One of ABY3 great strengths is its composability.
Each operation done on one or more tables produces as output also a table, which is a valid input for another query. This allows to build larger complex applications out of many small ones, very similar like one would do with a pipes-and-filters architecture. 
Furthermore, ABY3 comes with a description of how aggregate functions like MAX, SUM, COUNT can be realized when utilizing ABY3. For example, the maximum operator can be evaluated with a recursive algorithm that computes the maximum of the first and second half of the rows. In theory, ABY3 is able to compute any polynomial time function of a table, in practice, the efficiency may differ between functions and may not always be sufficient.   For executing its MPC operations ABY3 relies mainly on secret sharing.


\paragraph{Prototype Implementations}
ABY3 demonstrates its capability in two prototype applications. One of them could be used by the states of the United States to help ensure the validity of voter registration records. In the United States, each state maintains its own list of registered voters. Through the highly sensitive nature of these records coordination between states to ensure their faultlessness is not trivial. For that reason, one person moving from one state to another may often result in being registered in both states, which would allow them to illegitimacy cast a vote in both of these states.
ABY3 demonstrates how it could be used by the states to detect such double registration while preserving the confidentially of the records. 



\subsection{underlying MPC technology}
ABY3 works within a 3 party setting. This is a conscious decision as the two partiy and tree party setting each provide their own advantages and disadvantages.
The third party allows to deploy more efficient algorithms that could not be deployed in a two-party setting. For example, oblivious permutations can be done in O(n) in a three-party setting instead of O(n log n) in a two-party setting. On the other hand, there are already established solutions for many problems in the two-party setting that are not easily extendable to a three-party setting. ABY3 guarantees security against a semi-honest threshold adversary that is bound by an honest majority. For executing its MPC operations ABY3 relies mainly on secret sharing. As secret sharing comes with the advantage that algorithms based on secret sharing can have their input present in secret shared form and their output also is secret shared. When considering composability this is a great advantage as having input and output in the same format, as it allows to directly feed the output of one operation as input into the next one. While other MPC techniques like oblivious transfer require either input or output to be in the clear and would need to expansively transform it after each operation. 
ABY3s key feature are its new protocols for joins based on a MPC based cuckoo hash table. With these new protocols it is possible to join n rows with only O(n) overhead.

% all protocols constant rounds of communication \\
%- O(n) overhead in for join where n is number of rows \\

%-unique join keys
%- does feature a LAN VS WAN comparison in benchmarks




\paragraph{Computing Joins}
One key task for computing any kind of join is identifying which rows have identical join keys. More precisely if for two tables X,Y and key columns $X_1$, $Y_1$ and any given i  their exists j such that $X_1[i]$ = $Y_1[y]$. Where X[i] denotes the i-th row of table X and $ X_1[i] $ the i-th entry of the first column of table X. 

ABY3 implements an algorithm that solves this problem using a secure cuckoo hash table T with two hash functions $ h_1 $ and $h_2$. 
In a first step each row of Y is inserted into the hash table, such that Y[i] is inserted into $ T[h_0(Y_1 [i]  )] $ or $T[h_1(Y_1 [i]  )] $. 
If $X_1[i]$ has a matching join key, such that  $X_1[i]$ = $Y_1[y]$,  the matching row can only be located in $ T[h_0(X_1[i])] $ or  $T[h_1( X_1[i])] $. 
Therefore in a second step a match can found by comparing $X_1[i]$ and $ T[h_0(X[i])] $,  $T[h_1(X[i])] $ in a secure way. The key challenge in this algorithm is the construction and usage of a secure cuckoo hash table that does not leak sensitive information. ABY3 implements such a hash table based on an oblivious switching network.
\paragraph{Oblivious Switching Network}
TODO









\
%In \cite{10.1145/3372297.3423358} ..



%Similar to the SQL-union the aby3-union does row-concatenation of two tables.  \todo[fancyline]{ todo :union und andere oparatioren die von standart sql abweichen beschreiben}Note that the semantic of the set union operator in aby3 does not 100 percent matches the semantic of the SQL union operator. 
\section{smcql}

\subsection{functionality}
Just like Prio+ and contrary to conclave and aby3, smcql is not a framework for developing MPC solutions, but rather an already complete MPC system. Unlike in the case if Prio+ the lack of flexibility that come with this is not problem. As the intended use for smcql perfectly fits our requirements. Smcql ...

\subsection{underlying MPC technology}


\subsection{querie optimization   }





- prototype for 2 parties can be expanded for more parties
- honest adversary
\section{Discussion of Frameworks}
In this section we discuss important properties of the frameworks and compare them. 
At the end of the section we provide a \hyperref[Summary]{table} that summaries our discussion.
In their goals and general approach Conclave and SMCQL are very similar. Both Conclave and SMCQL employ an existing MPC frameworks in a new way, with the primary goal to make MPC more accessible. Conclave even describe SMCQL as ``the most similar state-of-the-art system'' \cite{10.1145/3302424.3303982}. ABY3 is new framework, that devolved new MPC protocols, with the primary goal to be as much efficient as possible. 
Conclave, ABY3 and SMCQL all assume a passive adversary that can corrupt at most one party. Therefore their security assumptions and guarantees are quite comparable. 
\paragraph{Features} 
When comparing the implemented features of the frameworks, there is no clear best framework. Since ABY3 focuses its efforts on joins, ABY3 provides the largest selection of join operations. In particular, ABY3 is the only framework that supports left and right joins. One category of features that ABY3 completely lacks are pre-implemented aggregate functions. Conclave and SMCQL both support sum aggregation. In addition, SMCQL is the only framework that supports pre-implemented maximums and count operators. Therefore SMCQL is the framework which supports the most aggregate functions. It is notable that the only two features that implemented in every frameworks are the Equi-join and UNION operators.  
\paragraph{Composability}
One notable aspect of each frameworks is their approach towards composability. Both Conclave and ABY3 lay huge emphasis on it and ensure all of their operators are composable. During our work with the frameworks this mechanism has saved us a lot of time and effort, because it significantly simplified the implementation of our more complex use cases.
SMCQL does not achieve full composability. On example where SMCQL fails to archive composability are select statements. As the majority of SMCQL's select statements cannot be subject to another select statement afterwards. During our work with SMCQL its limited composability has proven to impractical. 
\paragraph{Documentation}
It is not a trivial task to compare documentations objectively. There are hardly any measurable metrics that provide information about the quality of a documentation. However, one such metric is the amount of comments in the source code. We have found that SMCQL has the least amount of comments, in SMCQL there is about one comment for every 20 lines of code. ABY3 contains significantly more comments, the code of ABY3 contains one comment for every 10 lines of code. Conclave contains the most comments because Conclave contains one comment for every 5 lines of code. This measurement is consistent with the experience we have gained while working with the frameworks. In particular, we have observed Conclave's comments to be significantly more useful than those of ABY3. %As Conclave is the only framework that features external documentation is should be undisputed that Conclave has the best external documentation.  
We also perceived Conclave quick start guide to be more helpful then the quick start guide of ABY3. We suspect the reason for this is the different approach of the two frameworks towards the quick start guide. Conclave demonstrates in its quick start guide how to perform a single complete calculation. ABY3 on the other hand demonstrates multiple different small functions.  

\paragraph{Experimental Software}
One shared property of ABY3, Conclave and SMCQL, is their implementations are experimental software that primarily serves a proof of concept. This manifest especially as the existence of numerous bugs and undesirable behaviour that are part of every framework. Conclave for example features a failure detector that is supposed, to detect, if a party crashed because of an  internal error and than abort the computation. This failure detector itself is extremely error-prone and tends to produce a significant amount of false positives, i.e. it aborts the execution of the protocol, despite no party actually being crashed. Such a false positive may for example be caused, by a party that is fully occupied with a computation and does not respond to messages of the failure detector in time. For some queries this problem is more significant then for others, in extreme cases we observed more then 50\% of all executions being aborted without justification. %Another example is SMCQL's exception handling. SMCQL  



Installing a secure multiparty computation framework and doing its install setup is often a complex and non-trivial task. To quote Hastings et al. ``Installing and running these systems can be challenging'' \cite{hastings2019sok} 
During our work we have observed this statement to be true for each of our candidates.
Especially for Conclave and SMCQL the instructions are often vague or incomplete. Conclave for example requires the installation of JIFF \cite{JIFF} a JavaScript library, as one of its dependencies. In order for JIFF to be usable for Conclave, one must actively chose to ignore several recommendations during its installation. The task is additionally complicated by the fact that, neither Conclave nor SMCQL provide sufficient test to locate errors in their setup precisely. For installing and doing its initial setup, ABY3 comes with a script that automates this task. In theory this approach simplifies the installation a lot. In practise the automated installation comes with its own issues. If, like in our case, such an installation does not function correctly, fixing its errors is a time consuming task, as it's lacks documentation and therefore it is a none trivial task to identify the root causes of it errors.


\begin{table}[H]
	\phantomsection
	\label{Summary}	
	\centering
	\caption{Summary of important properties of the framework}
	\begin{tabular}{|l|l|l|l|} 
		\hline
		& ABY3                 & Conclave             & SMCQL                 \\ 
		\hline
		2 Parties            & No                   & Yes                  & Yes                   \\ 
		\hline
		3 Parties            & Yes                  & Yes                  & No                    \\ 
		\hline
		Passive Adversary    & Yes                  & Yes                  & Yes                   \\ 
		\hline
		Active~ Adversary    & No                   & No                   & No                    \\ 
		\hline
		Composable           & Yes                  & Yes                  & No                    \\ 
		\hline
		Quick Start Guide    & Yes                  & Yes                  & No                    \\ 
		\hline
		UNION Operator       & Yes                  & Yes                  & Yes                   \\ 
		\hline
		Equi-join            & Yes                  & Yes                  & Yes                   \\ 
		\hline
		Left Join            & Yes                  & No                   & No                    \\ 
		\hline
		Right Join           & Yes                  & No                   & No                    \\ 
		\hline
		Theta Join           & No                   & No                   & Yes                   \\ 
		\hline
		SUM aggregate        & No                   & Yes                  & Yes                   \\ 
		\hline
		MIN/MAX aggregate    & No                   & No                   & Yes                   \\ 
		\hline
		Projection           & No                   & Yes                  & Yes                   \\ 
		\hline
		Automated Setup           & Yes                   & No                  & No                   \\ 
		\hline
		
		\multicolumn{1}{l}{} & \multicolumn{1}{l}{} & \multicolumn{1}{l}{} & \multicolumn{1}{l}{} 
	\end{tabular}
\end{table}

\section{Rejected Frameworks}
\paragraph{CipherCompute}
One candidate for our study was CipherCompute \cite{Cosmian}. With the CypherCompute framework it is possible to solve a huge range of MPC problems using Rust. These include SQL operations like joins that are of interest for us. Furthermore, CypherCompute provides a rich documentation, consisting of a full quickstart guide and several well documented example projects. CypherCompute utilises the SCALE-MAMBA \cite{aly2021scale} framework for its underlying MPC operations. SCALE-MAMBA  itself has evolved out of the well-known SPDZ \cite{SPDZ} protocol. Unfortunately the early access version of CypherCompute is out of maintenance by the time we conducting this study. Therefore we have decided to not include CypherCompute in our study.

\paragraph{Prio+}
Prio+ \cite{cryptoeprint:2021:576} is the next generation of the Prio \cite{201553} framework. Prio+ strives to maintain the same use and security as Prio, while significantly increasing performance compared to its predecessor. Prio Plus allows an arbitrary number of parties to jointly compute aggregated statistics, like SUM, MAX, MIN operators. Prio+ utilises a client server model in which the (potentially many) input parties use a small number of servers to compute the statistics. Prio+ guarantees confidentiality of the input values if at least one server stays honest. Unlike CipherCompute or Conclave, Prio+ is not a framework for developing MPC solutions. Its rather a system for special purposes. This means that the use of Prio+ can not be extended beyond the usecases that have been originally implemented by the authors of Prio+. This leaves Prio+ with a relatively small range of usecases compared to frameworks like ABY3 or Conclave. Therefore we have decided to not include Prio+ in our study.

\paragraph{VaultDB} 
VaultDB \cite{rogers2022vaultdb} is a secure multiparty computation framework that was released in 2022. VaultDB is secure in the presence of a passive adversary and supports two parties, of which the adversary can corrupt at most one, therefore it resides in a setting with an honest majority. VaultDB is implemented on top of the EMP toolkit EMP toolkit \cite{emp-toolkit} which is an open source multiparty computation framework released in 2016. VaultDB provides an open source implementation which demonstrates its core capabilities.
A major shortcoming of VaultDB is the complete lack of documentation, and as a result, a very low usability. Therefore it has been proven to be infeasible for us to implement any none trivial use case within our time and resources. Consequently, it is not included in our study.

 
	
\chapter{Benchmarking}
\phantomsection
\label{Benchmarking}
For benchmarking efficiency there is often a variety of different metrics that are relevant and require to be measured with great care, as flawed or unclean measuring may deteriorate the value of the results. In this chapter, we describe the different metrics we want to benchmark and the different tools we use to achieve valid results. 

\section{Measuring Runtime}
A basic metric of how well a program works, is its runtime. Time is often measured in either wall-clock time or process time. Wall-clock time references, as the name implies, the passing of time on a wall-clock while the program runs. 
Process time resembles the actual time a CPU was used by the program. If, for example, the process blocks for a longer period of time, its wall-clock and process execution time may differ by a large amount. Typical reasons for a for a process to block would be that it waits for a slow hard drive or for incoming network traffic from another party. One needs to be careful what time is measured and that it is measured precisely. Otherwise one may obtain flawed or unfairly biased results. We measure both wall-clock and process execution time, as this allows us a far a more detailed analysis. In particular, this approach enables us to evaluate the difference between process time and wall-clock time, which indicates how long a process was blocked. In order to do so, tool-aided measuring is required. Measuring the individual frameworks is always trivial. Since the frameworks were implemented in different languages, different tools are available for measuring the individual frameworks. 
  \phantomsection
\label{time}	
\paragraph{SMCQL}%TODO blackbox metapher.
\code{time} \cite{time_sh} is a command line interface(CLI) command that is designed to benchmark the execution of other CLI commands. It is feasible to wrap the execution of a secure multiparty protocol in a .sh script and measure the performance of the script. Therefore \code{time} is in theory able to measure any of frameworks. One of \code{time's} advantages are the many different metrics it is able to measure, these include wall-clock and process time but also average memory requirements and maximum memory requirements. One major disadvantage of \code{time} is its inability to measure individual parts of a protocol in more detail. For SMCQL \code{time} perfectly fits our requirements. Therefore we use \code{time} to benchmark SMCQL  


\paragraph{Conclave}
Because \code{time} is able to measure shell commands we could also use it measure Conclave. Conclave is based on Python and Python comes with a ``batteries included'' approach, as it features a comprehensive standard library. Therefore plenty of tools, that provide valuable utility for us, are already integrated within python, because of this, there are other tools that are also valid candidates to measure Conclave. Two such candidates are timeit \cite{time} and the python profiler \cite{cProfile}, both are python libraries that offer a simple way to measure wall-clock or process execution time. Timeit measures exclusively end-to-end execution time, or in other words, the time the execution takes from one end to another. The python profiler comes with a more detailed analysis that includes detailed information on which functions have been executed, how often they have been executed, and how long it took to execute them. The extra utility provided by the python profiler does not come for free, compared to timeit, it has significant overhead that slows the execution down \cite{cProfile}. Therefore using it would result in an unfair disadvantage for Conclave. One significant disadvantage of timeit, is its inability to measure wall-clock and process time simultaneously. For this reason, in order to measure both of these values, it is required to perform each measurement twice, which doubles the time required and is impractical. Therefore we use the \code{time} tool with which we also measured SMCQL's execution time.  



\paragraph{ABY3}
 To benchmark ABY3 we use the cryptoTools library \cite{cryotoTools}, which is integrated into ABY3. The cryptoTools library is a C++ toolbelt that features a variety of tools for building cryptographic protocols. Among these utilities is a benchmarking tool for measuring runtime. With cryptoTools it is possible to measure the execution time of specific parts of the protocol. The croytoTools library provides the features we need and is already ``inbuilt'' into ABY3, therefore we have decided to use it for measuring ABY3's runtime. For a detailed example of how we utilise this tool benchmarking, see the description of our implementation for the second use case in \hyperref[Implementation]{Section 5}.




\section{Networking}
In our standard setup all parties run on the same machine and communicate through localhost. This simulates a practically perfect local area network (LAN) connection with very low latency and high throughput. It is also of interest how well the frameworks works in less ideal conditions. In particular, we also simulate a suboptimal wide area network(WAN) connection with high latency and limited bandwidth.

In order to do so we require a proxy server. Instead of connecting the parties to one another we connect them to the proxy server and the proxy server forwards the incoming messages to the addressed parties. To simulate a slow connection with high latency the proxy server delays incoming messages. Setting up such a proxy server is non trivial task, one challenge in particular is mapping the various connections to one another, in a correct way. This task is made more difficult by the fact that the different frameworks handle their connections in various different ways. In ABY3 for example, each party holds one direct connection to every other party. Differently, in Conclave, every party is connected to a Node.js server that forwards the messages. In order to handle these various approaches correctly, an analysis of the communication patterns is required. An additional factor that complicates our task is the fact that different frameworks use various different protocols to communicate. For example Conclave uses among others HTTP,while ABY3 utilities plain TCP. Checking implementation details and source code is a target-oriented approach for such an analysis , another tool that helped us to understand the communication patterns is Wireshark.      
\paragraph{Wireshark}
Wireshark \cite{wireshark} is an open source packet sniffer that allows to capture network traffic and saves it for a detailed analysis. Wireshark supports the analysis of numerous different protocols, among these are plain TCP, HTTP and websockets. Hence Wireshark supports all the protocols that our frameworks relies on, and therefore are of relevance for our work. With Wireshark we have been able to record the communication of our parties and pin down the exact communication patterns. Another utility Wireshark provides for us, is the ability to record communication once our proxy was setup. We also used Wireshark's recordings to confirm that our proxy works correctly once it was setup and used it detect error in our configuration. 
                

 \paragraph{Toxiproxy}
 Both ABY3 and SMCQL implement communication between parties based on a plain standard socket. In the case of ABY3 it is the standard C++ socket. For SMCQL it is the standard java socket. Both of these are TCP based and can be proxied with a standard TCP proxy. For this purpose we use Shopify's Toxiproxy \cite{toxiproxy}. Toxiproxy is a Go framework that allows to simulate different hazardous network conditions. These include a connection that delays its messages to simulate a high latency setup. Once the proxy server is setup it can be configured over the command line interface (CLI) or alternatively over an HTTP interface, Toxiproxy provides multiple different dedicated HTTP clients for this purpose. The clients differ in that they offer an interface in different programming languages. We have chosen to use the provided Ruby client as it is the one that provides the most extensive documentation. 
 A simplified example on how to use Toxiproxy to simulate latency can be found in \hyperref[the_label]{Listing 5.1}. In order to use Toxiproxy, we first set up the proxy so it starts to accept new connections. This is done by calling Toxiproxy populate and specifying to which address the proxy listens to and the address the messages get forwarded to.
 By default Toxiproxy does not add any network limitations to a connection. In our example we apply two limitations, which simulate a latency of 1000ms. The first limitation is applied to the ``upstream'' direction. Therefore it affects every message that is send towards the server from the address the server listens to. The second limitation is applied to the ``downstream'' direction. Therefore it affects every response that comes from the address the server listens to.    
 
   \phantomsection
 \label{Toxi_label}				
 \begin{lstlisting}[caption={Setting up a proxy that simulates latency between two parties with Toxiporxy}]
 #First we instantiate a connection between the two parties. 
 Toxiproxy.populate([
 {
 	name: "aby3_party2_party1",
 	#party 3 sends its messages for party1 to port 50010 therefore the proxy must listen to this port
 	listen: "127.0.0.1:50010",   
 	#party 1 listens to port 50001 therefore theproxy must forward to this port
 	upstream: "127.0.0.1:50001"  
}
 ])
 #Then we simulate a latency of 1000ms 
 toxiproxy-cli add aby3_party2_party1 -type latency -name upstream latency=1000 -upstream
 toxiproxy-cli add aby3_party2_party1 -type latency -name downstream latency=1000 -downstream
 
 \end{lstlisting}
 \paragraph{Node-Http-Proxy}
 Conclave's communication is based partially on websockets and partially on plain standard HTTP.
 Websockets are implemented on top of TCP. In particular, websockets use a single TCP socket for bidirectional communication. Therefore proxying Conclave cannot be done with a simple TCP proxy. Instead we use node-http-proxy a library for proxying HTTP that also supports websockets. With node-http-proxy we are able to delay messages to simulate high latency. It is also possible to measure the amount of data sent and received over a connection.
 Node-http-proxy is based on JavaScript and relies heavily on an event driven programming paradigm. A simplified example how to use node-http-proxy can be found in \hyperref[the_label]{Listing 5.2}. In this example we first create a proxy server and specify the address it forwards to. In a second step we create a a HTTP server that delays every incoming message for 500 milliseconds and then sends it to the proxy server. Subsequently we demonstrate how Node-Http-Proxy make use of the event driven programming paradigm. In the third step we subscribe to the ``proxyReqWs'' event. The ``proxyReqWs'' event is raised by the proxy each time an outgoing websocket message is received. More concrete event is raised before the message is transmitted to its destination and we can submit an anonymous function as event handler. The event handler is executed for each message before the message is transmitted. This for example enables to modify responses or to apply various other practical use cases. In our specific example, we provide a function that waits 500 milliseconds to simulate latency. In the fourth step we subscribe to the ``close'' event. The ``close'' event is raised each time a websocket is closed. We use our event handler to save the amount of bytes transmitted though the websocket. This is a showcase example, of how a proxy helps to measure important metrics that would be non trivial to measure otherwise.
  \phantomsection
  \label{the_label}				
 \begin{lstlisting}[caption={Setting up a proxy that simulates latency with node-http-proxy}]
	#create proxy server and specify the adress it forwards to
	#with the ws:true parameter we enable support for websockets
 	var proxy = new httpProxy.createProxyServer({
 		target: {
 			host: 'localhost',
 			port: 9005
 		},
 		ws:true 
 	});
 	 # Here we crate a standart HTTP server that delays every incoming message for 500ms and then forwars it the proxy server.
 	var proxyServer = http.createServer(function (req, res) {
 		setTimeout(function () {
 			proxy.web(req, res);
 		}, 500);
 	}).listen(9000);
 	# for every outgoing message the proxy emits a proxyReqWs event that we react to and delays the message for 500 milliseconds
 	proxy.on('proxyReqWs', function(){
 		 setTimeout(function(){ },500) 
 	 });
  	# each time a connection is closed the proxy emits a "close" event that we react to and save the amount of bytes transmitted through the connection
 	proxy.on('close' ,function (res,socket,head) { 
 		send  = socket.bytesRead;
 		received  = socket.bytesWritten;
 	});
 \end{lstlisting}

\section{Measuring Memory Consumption}
Another metric we collect is the consumption of memory. Sometimes it is possible to speed up computations at the cost of their memory consumption, so it is not sufficient to measure only runtime to get a fair measure of efficiency. A well-known example of such a time-memory trade-off are Rainbow Tables \cite{rainbow}. Another reason to measure memory usage is the fact that during our survey it often turned out to be a limiting factor for the size of datasets that are feasible to evaluate. 
For each framework we have measured its maximum memory consumption during an execution. With \hyperref[time]{time} we already have a tool that is able to measure the memory consumption of each framework in a sufficient way, therefore we use \code{time} to do so.


 
\chapter{Implementation}
In this chapter we describe the use-cases we have chosen to implement and benchmark. We first describe and motivate our choice of use-cases and, afterwards, give reasonable details on their implementation. We do implement four use-cases that are ordered from the least complex to the most complex.
\section{Use-Case Description} 
We have decided to implement every use-case with Conclave and SMCQL for two parties. As ABY3 does not allow a two party protocol but requires at least three parties, we have implemented everything for ABY3 with three parties. For ABY3 there will always be only two parties that provide input data. The third party will not provide input data but will assist in the execution of the protocol. We consider this approche justified, as it is the exact approach Conclave and ABY3 use in their Evaluation. And if the authors of the frameworks considers such an approach fair, then we do hold no objections against it.  
We will refer to first party that provides input as Alice and to the second party that provides input as Bob. Furthermore if a framework is not capable of replicating the semantic of a use-case to one hundred percent, we do use that functionality available to implement the most similar semantic the framework can provide. For example use-case two features boolean logic, but Conclave does not feature an explicit boolean logic. Therefore we use integers where the integer one represents the value true and the integer zero repents false.   
\paragraph{Use-Case One}
For our first use-case, we have chosen a single join. A single join may not be a very complex use-case but it is not irrelevant as joins are of great importance for practically every relational database query. It is estimated that over 60 \% of privacy-sensitive analytics queries include at least one join \cite{johnson2017practical}. In our first experiment, Alice and Bob each hold one table. Each of these tables consists of 4 columns. The first column serves as the primary key that is used for the join. The other 3 columns are filled with random non-negative integers and simulate user data. We have chosen to use non-negative integers, because the usage of negative integers has resulted in an  increased frequency of various bugs and other undesirable behaviour. We are calculating an \hyperref[Databases]{equijoin}  with the primary key generated in such a way that 50 \% of the entries in each table will match the join criteria. We do not generate the primary key randomly, because they size of output relation depends on it and we want those sizes to be consistent between individual executions. As result, we will reveal the entire outcome of the join. So the primary utility provided by the use of MPC operations is the fact that the entries that do not match the join criteria are obscured.

  \phantomsection
\label{SQL1_label}				
\begin{lstlisting}[caption={ Functional equivalent SQL statement for our first use-case  }]
SELECT * 
FROM Alice A JOIN Bob B 
ON  A.primary_key = B.primary_key
\end{lstlisting}
   \phantomsection
\label{SQL1}
			
\paragraph{Use-Case Two}
Computing joins alone is of limited use if the result of the join can not be subject to further selection. Therefore in our second use-case we will first compute a join and the query the result with a classic "SELECT ... FROM ... WHERE" statement. Similar to first experiment Alice and Bob again each hold one table. The tables consist of two columns. The first column serves as primary key and the second column contains a boolean values that is generated at random. For our experiment we will in a first step compute the natural join of the two tables, where the primary key column zips the two tables together. The primary keys will be again generated in such a way that 50 \% of the columns will be included in the join. In a second step we will apply a where filter to the result of the join and eliminate every row that does not have two identical boolean values. This use case functions also as a simple showcase example for composability and will show how well this mechanism functions in practise.    
\begin{lstlisting}[caption={Functional equivalent SQL statement for our second use-case}]
	SELECT PrimaryKey, AliceBool, BobBool
	From AliceTable 
	NATURAL JOIN BOB TABLE
	WHERE AliceBool = BobBool
\end{lstlisting}
\paragraph{Use-Case Three}
Besides joins, another very important group of SQL operations are aggregate functions. Over a third of all privacy-sensitive analytics queries requires a aggregation \cite{johnson2017practical}. Therefore our third use-case is centred around an aggregate function, or more precisely a maximum operator.  TODO
%In our third experiment Alice and Bob each hold one table. The tables consist of two columns each. Like in our first and second use-case the first column will function as primary key and will be used for a join. The second column will be the column we aggregate over. In a first step we will be compute the join of the two columns  
 

\paragraph{Use-Case Four}
For our fourth and last use-case, we compare two special features of SMCQL and Conclave. Both Conclave and SMCQL feature a mechanic, that allows revealing some of the columns of the input data. The revelation of the input data allows them to apply optimizations that speed up computation while preserving the privacy of the other columns. For a more detailed description see \hyperref[Accesses_label]{SMCQL's access control} and \hyperref[Trust_label]{Conclave's trust annotations}. Therefore in our fourth use-case, we are going to replicate the setup of our first use-case but this time we will allow the leakage of the primary key column. Such leakage will allow both Conclave and SMCQL to apply their optimization. Replicating the setup of the first use-case enables us, to compare the results of the fourth use-case to the first use-case. This comparison will show how big the speed up of these optimizations is in practice. 

















\section{Implementation}
   \phantomsection
\label{Implemetation}	
In this section we provide a description of notable details of the implementation of our use cases. 
The descriptions are ordered from the least complex to the most complex use case. Since SMCQL only requires the query to be specified in a valid SQL syntax and does the rest of the implementation automatically, we focus here on Conclave and ABY3 as their implementation is more complex.    
\paragraph{Use Case One}
As our first use case is the least complex one, we use it for a full demonstration how ABY3's and Conclave's basic workflow functions. 
\paragraph{Conclave} 
The first step in Conclave is always the definition of the relation scheme of the input. Accordingly, we first define two lists of columns, one for Alice and one for Bob. Each column is defined with a name, a data type, and an integer that indicates its owner. The owner is used to annotate trust and apply optimizations, as described in our framework description of \hyperref[Trust_label]{Conclave}. In our first use case each party trust itself and only itself. Once the relation scheme is defined we can populate it by using the \code{create} function. Create loads the data from a .CSV(Comma-Separated-Values ) file that we generated previously. Once the tables are populated we can compute the join. In order to, compute a join in conclave we need to parse the two tables of the join, a name for the result, and the two columns that are used as key columns for the join. 
Lastly, we can reveal the output of our computation using the \code{collect} function. The \code{collect} function requires two parameters, one that specifies which table is revealed and a second one that specifies to whom it is revealed.
\begin{lstlisting}[caption={ The Python protocol of Conclave for our first use case    }]
def protocol():
	# define the schema for the input tables 
	AliceColumns = [
	defCol("primary_key", "INTEGER", [1]),
	defCol("user_data1_Alice", "INTEGER", [1]),
	defCol("user_data2_Alice", "INTEGER", [1]),
	defCol("user_data3_Alice", "INTEGER", [1]),
	]
	BobColumns = [
	defCol("primary_key", "INTEGER", [2]),
	defCol("user_data1_Bob", "INTEGER", [2]),
	defCol("user_data2_Bob", "INTEGER", [2]),
	defCol("user_data3_Bob", "INTEGER", [2])
	]
	# the content of the tables is loaded from a pregenerated .CSV
	AliceTable = create("AliceTable", AliceColumns, {1})
	BobTable = create("BobTable", BobClumns, {2})
	# calculate the join over the two tables 
	JOIN = join(AliceTable, BobTable, 'JOIN', ['primary_key'], ['primary_key'])
	# reveal the output of the join to Alice
	collect(JOIN, 1)
	# reveal the output of the join to Bob
	collect(JOIN, 2)
\end{lstlisting}

\paragraph{ABY3}
Before the actual execution of the protocol starts ABY3 requires some setup. Similar to Conclave, the first step in ABY3 is also the definition of the relation scheme. Hence we start with the definition of two lists of \cite{ColumnInfos}. Each column is described by a name, a data type, and some data type-dependent information. In our case, we use integers and annotate that we use 32-bit integers. After the relation scheme is defined we create local tables and populate them. One last step that needs to be done before the execution of the protocol starts is the initialization of a message server that manages the communication between the parties.
Once the execution of the protocol starts we need to convert our local input into a table that is secret shared between the parties. 
Before we can compute the join of the shared tables we need to define which columns of the shared tables we want to select. Therefore we define a list of references to the shared columns. In order to to obtain the desired semantic, the list needs to contain each unique column and one reference to each column that is a duplicate in both tables. 
In our use case the columns containing the user data are unique and the key columns are duplicates because they have identical naming and data types. Therefore our list contains seven columns. In the next step we compute the join by providing the join function with one reference to each key column and a list of columns we want to select. The output of the join is still a shared table. To obtain the plaintext values of our result we need to explicitly convert the shared table.  


\begin{lstlisting}[caption={Simpifiyed Protocol for our first use case in ABY3}]
void use_case_two(u32 rows, u32 cols ){
	std::vector<ColumnInfo> AliceCols = { ColumnInfo{ "key", TypeID::IntID, 32}};
	std::vector<ColumnInfo> BobCols = { ColumnInfo{ "key", TypeID::IntID, 32}};
	for (u32 i = 1; i < cols; ++i)
	{
		AliceCols.emplace_back("Alice" + std::to_string(i), TypeID::IntID, 32);
		BobCols.emplace_back("Bob" + std::to_string(i), TypeID::IntID, 32);
	}
	# Create tables for Alice and Bob and fill them with content
	Table AliceTable(rows, AliceCols);
	Table BobTable(rows, BobCols);
	u32 intersectionsize = rows * 0.5;
	# Fill the primary columns such that 50% of all entries match the join criteria 
	for (u64 i = 0; i < rows; ++i)
	{
		# if out is false then the entry will be included in the join
		auto out = (i >= intersectionsize);
		AliceTable.mColumns[0].mData(i,0) = i + 1;
		BobTable.mColumns[0].mData(i, 0) = i + 1 + (rows * out);
	}
	# Fill the other columns with random integers
	for (u64 i = 1;  i < cols; ++i){
		for (u64 j =0; j < rows; ++j){
			AliceTable.mColumns[i].mData(j, 0) = rand() ;
			BobTable.mColumns[i].mData(j, 0) = rand();	
		}
	}
	# Instanciate server that handels communication
	DBServer server[3]; 
	...
	# Here the exectution of the protocoll starts.
	auto protocoll = [&](int i) {
		# create a secret shared verion of the input
		SharedTable AliceSharedTable = (i == 0) ? 
		server[i].localInput(AliceTable) : server[i].remoteInput(0);
		SharedTable BobSharedTable = (i == 1) ?
		server[i].localInput(BobTable) : server[i].remoteInput(1);
		# define a list of clumns we want to select
		std::vector<SharedTable::ColRef>  Select_columns;		
		for (u64 i = 0;  i < cols; ++i){
			Select_columns.emplace_back(SharedTable::ColRef(BobSharedTable,BobSharedTable.mColumns[i]));
		}
		for (u64 i = 1;  i < cols; ++i){
			Select_columns.emplace_back(SharedTable::ColRef(AliceSharedTable,AliceSharedTable.mColumns[i]));
		}
		# compute the join
		SharedTable JOIN =srvs[i].join( SharedTable::ColRef(AliceSharedTable,AliceSharedTable.mColumns[0]), SharedTable::ColRef(BobSharedTable,BobSharedTable.mColumns[0]), First_Select_columns);
		# reveal the plaintext values of the result
		aby3::i32Matrix result = RevealAll(server[i], JOIN); 
	};
	
	auto AliceThread = std::thread(protocol, 0); start Alice's thread
	auto BobThread = std::thread(protocol, 1); start Bob's thread
	thread(protocol, 2); # start the assisting third party
	t0.join(); # wait for Alice to finish
	t1.join(); # wait for Bob to finish
}

\end{lstlisting}
\paragraph{use case Two}
For our second use case, we do not implement the query directly. Instead, we apply an optimized query that has an identical output. To obtain the desired result Alice and Bob compute two intermediate results, the union of these intermediate results will then yield our final result, an equivalent SQL statement is listened \hyperref[use_case2_alternative_sql]{ in Listing 6.5}.
In the first step, Alice and Bob filter their input once so that it contains only entries with the value false in its boolean column. Afterward, they can compute the join of these two filtered tables. The result of this join is the first important intermediate result, as each of its entries is also part of the final result. The second important intermediate result can be obtained by the same procedure when filtering the input for entries that contain true.
The union of these two intermediate results is our final result. 

We have chosen to implement the use case in this indirect way for two reasons. First Conclave cannot evaluate the query directly, as Conclave's current prototype implementation can apply a WHERE filter exclusivity to a table that is either the direct input of a party or the output of a unary operator, that has only one table as input. Additionally, our indirect implementation has a significantly better performance compared to the direct implementation, for a detailed analysis see our evaluation in \hyperref[evaluation]{Chapter 7}.   


\phantomsection
\label{use_case2_alternative_sql}
\begin{lstlisting}[caption={Functional equivalent SQL statement for our optimitized implementation of our second use case}]
SELECT PimaryKey, AliceBool, BobBool FROM
(SELECT PrimaryKey, AliceBool FROM AliceTable WHERE AliceBool == false) 
NATURAL JOIN	
(SELECT PrimaryKey, BobBool FROM AliceTable WHERE BobBool == false)	
UNION
SELECT PimaryKey, AliceBool, BobBool FROM
(SELECT PrimaryKey, AliceBool FROM AliceTable WHERE AliceBool == true) 
NATURAL JOIN	
(SELECT PrimaryKey, BobBool FROM BobTable WHERE AliceBool == true)
\end{lstlisting}

\paragraph{Correctness}
Despite it not being obvious, our procedure indeed yields the correct result, as each join generates only entries, that have a primary key that is included in Bob's and Alice's inputs and have identical boolean values. Therefore every entry that is part of our result is also part of the correct result.

If any entry is not part of our result, it cannot be part of either intermediate result. Any entry that is not part of either intermediate result, features a primary key that is either not represented in both inputs or contains two boolean values that are not equal.
Therefore the entry cannot be part of the correct result. Consequently, each entry that is not part of our result, cannot be part of the correct result. Which is the contraposition and therefore equivalent, to the fact that each entry of the correct result is part of our result. 
As our result contains every entry of the correct result and the correct result contains every entry of our result, our result, and the correct result are identical. Therefore our procedure is indeed correct.

\paragraph{Conclave}
In order to implement our second use case in Conclave, we start with the import of the input. Therefore we annotate the layout of each table. Conclave does not support a dedicated data type for boolean values, consequently, we use integers for every column. In the second step, we generate four new tables by applying the filters to the input. These filters are a showcase example of how Conclave can speed up computation by applying optimizations.
A naive approach for applying the filters would be to share the input tables and then filter the shared tables using an MPC algorithm. This is an approach Conclave is in theory capable of. But Conclave is able to utilize the fact, that each filter has only a single input that is known to one party in the clear. In Consequence, the parties can apply the filter to their input locally without the use of an MPC algorithm. After they have applied the filters the parties can then create shared tables of the filtered input. In the next step compute the two intermediate results by joining the shared filter input of Alice and Bob. Here Conclave demonstrates how it supports composability, as the output of the filter function is used as input for the join function.  
Finally, we compute the row-wise concatenation of the two intermediate results and publish the final result to Alice and Bob.   
\begin{lstlisting}[caption={Simpifiyed Protocol for our second use case in Conclave}]
def protocol():
	AliceColumns = [
	defCol("primary_key", "INTEGER", [1]),
	defCol("AliceBool", "INTEGER", [1])
	]
	AliceTable = create("AliceTable", AliceColumns, {1})
	BobColumns = [
	defCol("primary_key", "INTEGER", [2]),
	defCol("BobBool", "INTEGER", [2])
	]
	BobTable = create("input_2", AliceColumns, {2})
	AliceFilterFalse = cc_filter(AliceTable, "AliceFilterFalse", "AliceBool", "==", scalar=0)
	BobFilterFalse = cc_filter(BobTable, "AliceFilterFalse", "BobBool", "==", scalar=0)
	AliceFilterTrue = cc_filter(AliceTable, "BobFilterTrue", "AliceBool", "==", scalar=1)
	BobFilterFalse = cc_filter(BobTable, "BobFilterTrue", "BobBool", "==", scalar=1)
	intermediateFalse = join(AliceFilterFalse, BobFilterFalse, 'intermediateFalse', ['primary_key'], ['primary_key'])
	intermedidateTrue = join(AliceFilterTrue, BobFilterTrue , 'join_result1', ['primary_key'], ['primary_key'])
	final_result = concat([ intermediateFalse , intermediateTure ], "final_result" )
	collect(final_result, 1)
	collect(final_result, 2)
\end{lstlisting}



\paragraph{ABY3}
Our second use case is also a showcase example of how we use ABY3's integrated benchmarking capability to sample valuable information. With the in ABY3 integrated timer, we are able to set \code{TimePoints} at arbitrary potions of the protocol. After the execution is finished, the timer returns how much time is passed before each individual \code{TimePoints} was reached. For this reason, we obtain detailed information, on how the different parts of the computation compose its total runtime. Since Conclave is able to apply the filter locally, in order to hold the comparison between them fair, we also filter the locally. Therefore Alice and Bob apply the filters to their input using standard C++ before the protocol begins. The majority of the work is done by two functions we present in Listing 6.6. The first Function generates the intermediate results, in order to do so, in the first step we load the input and convert it into a secret shared table. In the next step we calculate the join over the shared tables, to do so we need to pass a reference towards the two key columns, together with a list of the columns we want to select over. Our second function uses the first function to obtain the intermediate results. Once we have obtained those we compute their union.           
 \begin{lstlisting}[caption={Simpifiyed Protocol for our second use case in ABY3}]
	# include the timer and instanciate it
	include "cryptoTools/Common/Timer.h"
	Timer t;
	...
	# genarte input and apply filter to input locally, omited for reason of simplicty.
	...
	# i denotes the party id where i==0 means Alice, i==1 Bob and i==2 denotes the third assisting party
	auto getIntermediate =[&] (int i, int filter){
		SharedTable AliceTable;
		SharedTable BobTable;
		#  Load the prefilterd input and convert it into a secret share. 
		if (filter == 1) { 
			AliceTable = (i == 0) ? srvs[i].localInput(AliceInputFiltertOne): 
			srvs[i].remoteInput(0); 
			BobTable = (i == 1) ? srvs[i].localInput(BobInputFiltertOne): 
			srvs[i].remoteInput(1);
		}else{
			AliceTable = (i == 0) ? srvs[i].localInput(AliceInputFiltertZero): 
			srvs[i].remoteInput(0);
			BobTable = (i == 1) ? srvs[i].localInput(BobInputFiltertZero): 
			srvs[i].remoteInput(1);
		}# Here we annotate that we want to select the first two columns of Alice's Table and the second collumn of Bob's Table
		std::vector<SharedTable::ColRef>  First_Select_collumns; 
		Select_columns.emplace_back(SharedTable::ColRef(BobTable, BobTable.mColumns[0]));
		Select_columns.emplace_back(SharedTable::ColRef(BobTable, BobTable.mColumns[1]));
		Select_columns.emplace_back(SharedTable::ColRef(AliceTable, AliceTable.mColumns[1]));
		# We caluclate the join and return it
		return srvs[i].join( SharedTable::ColRef(AliceTable , AliceTable.mColumns[0]), SharedTable::ColRef( BobTable, BobTable.mColumns[0]), Select_columns);
		
	};
	auto use_case2 = [&](int i, int filter)
	{
		t.setTimePoint("start"); 
		SharedTable IntermediateOne = getIntermediate(i,0);
		t.setTimePoint("FirstIntermediate"); 
		SharedTable ItermediateZero = getIntermediate(i,1);
		t.setTimePoint("SecondIntermediate");
		SharedTable UNION;
		# Here we annotate which columns of the first Intermediate 
		# will part of the UNION
		std::vector<SharedTable::ColRef>  SelectOnes; 
		std::vector<SharedTable::ColRef>  SelectZeros; 
		for(u64 index=0; index <3; ++index)  {
			# Here we annotate that all columns of both intermediateresults 
			# will be part of the UNION
			SelectOnes.emplace_back(SharedTable::ColRef( IntermediateOne , 
			IntermediateOne.mColumns[index]));
			SelectZeros.emplace_back(SharedTable::ColRef(IntermediateZero , 
			IntermediateZero.mColumns[index]));
		}
		t.setTimePoint("UNION");
		# Compute the union of the two intermediate results
		UNION = srvs[i].rightUnion(SharedTable::ColRef( IntermediateZero, IntermediateZero.mColumns[0]),
		SharedTable::ColRef( intermediateOne, intermediateOne.mColumns[0]),  SelectZeros,SelectOnes);
		t.setTimePoint("end");
		
	};
	auto t0 = std::thread(use_case2, 0,1); #Start Alice thread
	auto t1 = std::thread(use_case2, 1,1); #Start Bob thread
	use_case2(2,1); # Start the assisting third party
	t0.join(); # Wait for Alice to finish
	t1.join(); # Wait for Bob to finish
	write_to_file(t) # Write timer information to file for further analysis
\end{lstlisting}

\paragraph{use case Three}
comming soon sollte aber nicht so viel sein da nicht mit aby3 implmentiert wird

\paragraph{use case Four}
As our first and fourth use case a nearly identical their implementation differs only in minor details. Especially for Conclave the implementation is identical besides the fact that we need to use different trust annotations. In order to allow Conclave to apply its public join optimization, we annotate that each party is allowed to learn both of the key columns.

\begin{lstlisting}[caption={ The Python protocol of Conclave for our last use case    }]
	def protocol():
	# define the schema for the input tables 
	AliceColumns = [
	# annotate that Alice trust Bob to learn hear key column
	defCol("primary_key", "INTEGER", [1,2]), #
	defCol("user_data1_Alice", "INTEGER", [1]),
	defCol("user_data2_Alice", "INTEGER", [1]),
	defCol("user_data3_Alice", "INTEGER", [1]),
	]
	BobColumns = [
	# annotate that Bob trust Alice to learn his key column
	defCol("primary_key", "INTEGER", [1,2]),
	defCol("user_data1_Bob", "INTEGER", [2]),
	defCol("user_data2_Bob", "INTEGER", [2]),
	defCol("user_data3_Bob", "INTEGER", [2])
	]
	...
	...
	collect(JOIN, 1)
	collect(JOIN, 2)
\end{lstlisting}




\chapter{Evaluation}
\phantomsection
\label{Evaluation}	
In this chapter we discuss the results of our benchmarks. We provide figures for the individual performance of the frameworks and highlight notable details. We evaluate the result of each use case individually, ordered from least complex to most complex. At the end of the chapter we discuss our results and provide our final conclusion.
\paragraph{Experimental Setup}
All our benchmarks are conducted on a single machine that hosts all parties. Our machine runs on a 64 bit Linux operations system, more precisely we use Debian 11 with the Linux Kernel version 5.10.0-18.amd64. The system features 4 CPUs that each run at 2,3 GHz and 128 GB of RAM. In order to fit the results of different frameworks into shared figures, we use a logarithmic scaling to base 10 for our figures. We always annotate the total amount of input-rows included in a computation. For example, if a computation involves two parties that each provide an input of 100 rows, we will annotate that as 200 rows of input. To obtain more reliable measurements, we perform each measurement multiple times, wherever possible, and select the median of our measurements as final result.
\section{Use case One}
For our first evaluation we are inspecting the result of our first use case in the localhost setting. Our first use case is a simple join for a detailed description see \hyperref[Use Case]{Section 5.1}. We have visualised a comparison of runtime and space consumption in Figure 6.1 and 6.2.

\paragraph{Required Runtime}
For SMCQL we have been able to evaluate the query for up to 140 input-rows. In order to compute the query with 140 input-rows SMCQL took about 10 hours. For Conclave we have been able to scale up to 500 input-rows, with 500 rows taking about 8 hours. Neither SMCQL nor Conclave could produce a result for larger datasets within 24 hours and after 24 hours we stopped the computation. A visitation of the result for our result can be found in Figure 6.1.
We have observed that the difference in speed between SMCQL and Conclave grows significantly with increasing input size. For an input of size 20, Conclave is 5 times faster than SMCQL, for an input of size 100 Conclave already is more than 50 times faster and for an input of size 140, Conclave is more than 75 times faster than SMCQL. In the first use case ABY3 is by far the best scaling framework. ABY3 is able to compute the query with 140 and 500 input columns in less than a second. ABY3 runtimes does not increase in any significant way before its input reaches the mark of 128000 input-rows, for which it requires 1107 milliseconds. We have been able to evaluate the query with 16.000.000 input-rows in one minute and 46 seconds. 
\paragraph{Evaluation of used Space}
It's notable that for large input sizes SMCQL is more space efficient than Conclave, while for small input input sizes Conclave is more space efficient than SMCQL. They break even at an input size of 60 rows, for which they both require about 870 MB of space.  For input larger than 60 rows SMCQL scales significantly better than Conclave. For an input of size 140 rows SMCQL only needs about 1300MB while Conclave needs more than 4 times as much. As the space consumption of SMCQL grows linear and the space consumption of Conclave doubles, roughly every time the input-size grows by 100, we can only assumes that this trend would continue for large datasets. Similar to the case of the runtime, ABY3 also performs significantly better than Conclave and SMCQL. For input size 140, ABY3 only allocates 28MB of space and therefore is over 46 times more efficient than SMCQL. For input size 500, ABY3 allocates about 32 MB of space and is therefore 1800 times more efficient than Conclave, which requires over 58 GB of space.          

\label{evaluation}
%10 6:17
%20 25:02
%40 2:55:29
%50 23:05
%60 7:24:52
%40 3:07:24
%50 4:00
%70 7:38
%100 17:24.14
%150 53:23
%200 2:57:41
%250 6:58:13 6.90775527898214
\begin{figure}[H]
\begin{tikzpicture}
		\begin{axis}[
		xmin = 0, xmax = 500,
		ymode=log,
		xtick distance = 50,
	%	log ticks with fixed point,
	%	x filter/.code=\pgfmathparse{#1 },
		grid = both,
		minor tick num = 1,
		major grid style = {lightgray},
		minor grid style = {lightgray!25},
		width = \textwidth,
		height = 0.5\textwidth,
		xlabel = {Size of Dataset},
		ylabel = {Runtime im ms},
		legend style={at={(0.5,-0.2)},
		anchor=north,legend columns=-1}   ]	
	
	\addplot[
	color=blue!50!cyan,smooth,tension=0.7,very thick
	] file[skip first] {conclave_use_case1_loca.dat};
	\addlegendentry{Conclave}
	
	%\addplot[color=blue!50!cyan,smooth,tension=0.7,very thick]file[skipfirst]{ABY_use_case1_local.dat};
	\addplot[
	color=orange!50!cyan,smooth,tension=0.7,very thick
	] file[skip first] {SMCQL_use_case1_local.dat};
	\addlegendentry{SMCQL}
	\addplot[
	color=green!50!red,smooth,tension=0.7,very thick
	] file[skip first] {ABY_use_case1_local.dat};
	\addlegendentry{ABY3}
	
	\end{axis}
	
\end{tikzpicture}
	\caption{Runtime of ABY3, Conclave and SMCQL for in our first use case}
\end{figure}
%1466320 + 1419280 + 335772
%2300692 + 2470452 + 780664
%4344428 + 4358604 + 1217576
%9497772 + 9710384 + 2649820 
%14496916 + 16821556 + 6677548
%25940748 + 26352440 + 5739388


\begin{filecontents}{space_conclave_1.dat}
	0 0
	0 0
	100 3221372
	140 5551808
	200 9920608
	300 21857976
	400 37996020
	500 58032576
\end{filecontents}
\begin{filecontents}{space_SMCQL_1.dat}
	0 0
	0 0
	50 3221372
	70 5551808
	100 9920608
	150 21857976
	200 37996020
	250 58032576
\end{filecontents}
\begin{figure}[H]
\begin{tikzpicture}
	\begin{axis}[
		xmin = 0, xmax = 500,
		ymode=log,
		xtick distance = 50,
	%	log ticks with fixed point,
	%	x filter/.code=\pgfmathparse{#1 },
		grid = both,
		minor tick num = 1,
		major grid style = {lightgray},
		minor grid style = {lightgray!25},
		width = \textwidth,
		height = 0.5\textwidth,
		xlabel = { Size of Dataset },
		ylabel = { Space consumption in kbytes},
		legend style={at={(0.5,-0.2)},
		anchor=north,legend columns=-1}]	
		
		\addplot[
		color=blue!50!cyan,smooth,tension=0.7,very thick
		] file[skip first] {space_conclave_1.dat};
		\addlegendentry{Conclave}
		\addplot[
		color=orange!50!cyan,smooth,tension=0.7,very thick
		] file[skip first] {space_SMCQL_1.dat};
		\addlegendentry{SMCQL}
		\addplot[
		color=green!50!red,smooth,tension=0.7,very thick
		] file[skip first] {space_ABY3_1.dat};
		\addlegendentry{ABY3}
	\end{axis}
\end{tikzpicture}
	\caption{Memory usage of ABY3, Conclave and SMCQL in our first use case}
\end{figure}
\section{Use Case Two} 
In \hyperref[use_case2]{Section 5.2} we have described an optimized way to evaluate the query of our second use case. In order to see if how large our efficiency gain is in practice, we have implemented this use case for ABY3 twice. One implementation is optimised to evaluate the query indirectly, such that in obtains a result with two joins of size n/2 and a local application of filters, the other implementation is fully unoptimized and obtains a result with one join of size n and  an application of filter with an MPC algorithm. Because our other result are very similar to the results of use case one, we focus on a comparison between the two implementation of ABY3.  TODO naive implementation im appendix wenn zeit
\paragraph{Evaluation of Runtime}
As Figure 5.3 depicts, our measurement shows that for small input sizes the unoptimised implementation is faster, while for larger input sizes the optimized implementation is faster. For a visualisation of our comparison see Figure 6.3. For input size 2000, the unoptimised implementation runs within 182 milliseconds, which is 1,5 times faster then our optimised implementation, that needs 276 milliseconds. With increasing input size, the difference between their performances becomes less significant, for input size 8000 the unoptimized implementation is less then 1,4 times faster. For an input of size 64000 the two implementation have very similar runtime and both need about 500 milliseconds. From there on onwards, our optimization starts to pay of and is constantly about 1,5 times faster then the unoptimised version. 
TODO ursache erklären warum erst langsamer dann schneller

\begin{figure}[H]
\begin{tikzpicture}
	\begin{axis}[
		xmin = 0, xmax = 16224000,
		ymin = 0, ymax = 200000,
		%xtick = {0,4000,8000,16000,32000,64000,128000,512000,1024000,2048000,4096000,8112000,16224000  },
		xmode= log,
		ymode= log,
		%ytick distance = 20000,
		grid = both,
		minor tick num = 2,
		major grid style = {lightgray},
		minor grid style = {lightgray!25},
		width = \textwidth,
		height = 0.5\textwidth,
		xlabel = { Size of Dataset },
		ylabel = {Runtime im ms },
		legend style={at={(0.5,-0.4)},
		anchor=south, legend columns=1}]
			
		
		\addplot[
		color=blue!50!black,smooth,tension=0.7,very thick
		] file[skip first] {ABY_use_case2_local.dat};
		\addlegendentry{Naive Implementation}
		\addplot[
		color=green!50!red,smooth,tension=0.7,very thick
		] file[skip first] {ABY_use_case2_fast_local.dat};
		\addlegendentry{Optimized Implementation}
	\end{axis}
	
\end{tikzpicture}
\caption{Runtime of our two implementations of ABY3  use case two}
\end{figure}
\paragraph{Evaluation of Space}
From the perspective of memory consumption our optimization is a strict improvement over the unoptimized implementation.
The optimized implementation is strictly better, as for every single input size we have tested, it requires less memory. We visualise the memory requirements of our two implementations in Figure 6.4.
It is notable how similar both implementations scale. As, for large inputs, they both very reliably double their memory consumption each time their input is doubled and the optimized implementation consistently takes half the memory of the unoptimised implementation. This trend start with input of size 64000 where our optimization needs about 300MB and the naive implementation about 150MB and continues from there on.
\begin{figure}[H]
	\begin{tikzpicture}
		\begin{axis}[
			xmin = 0, xmax = 16224000,
			ymin = 0, ymax = 63131608,
			%xtick = {0,4000,8000,16000,32000,64000,128000,512000,1024000,2048000,4096000,8112000,16224000  },
			xmode= log,
			ymode= log,
			%ytick distance = 20000,
			grid = both,
			minor tick num = 2,
			major grid style = {lightgray},
			minor grid style = {lightgray!25},
			width = \textwidth,
			height = 0.5\textwidth,
			xlabel = { Size of Dataset },
			ylabel = {Memory consumption in kb },
			legend style={at={(0.5,-0.4)},
				anchor=south, legend columns=1}]
			
			
			\addplot[
			color=blue!50!black,smooth,tension=0.7,very thick
			] file[skip first] {space_ABY3_2.dat};
			\addlegendentry{Naive Implementation}
			\addplot[
			color=green!50!red,smooth,tension=0.7,very thick
			] file[skip first] {space_ABY3_fast_2.dat};
			\addlegendentry{Optimized Implementation}
		\end{axis}
		
	\end{tikzpicture}
	\caption{Space requirement of our two implementations using ABY3 in use case two}
\end{figure}
\paragraph{Comparison between Conclave and ABY3}
In our second use case Conclave scales better than in the first one. For a visualisation of its results see Figure 6.5. We have been able to obtain a result for 1000 input columns in 6 hours and 32 minutes. Compared to the 500 input columns in 8 hours from the first use case, Conclave has been able to handle an input twice as large in less time. Yet despite these better results Conclave is still unable to compete with the results of our two ABY3 implementations for the second use case results, that both have been able to compute the result for 1000 input-rows in less then a second.

\begin{figure}[H]
	\begin{tikzpicture}
		\begin{axis}[
			xmin = 0, xmax = 1000,
			ymin = 0, ymax = 43099000,
			%xtick = {0,4000,8000,16000,32000,64000,128000,512000,1024000,2048000,4096000,8112000,16224000  },
			ymode= log,
			xtick distance=100 ,
			%ytick distance = 20000,
			grid = both,
			minor tick num = 2,
			major grid style = {lightgray},
			minor grid style = {lightgray!25},
			width = \textwidth,
			height = 0.5\textwidth,
			xlabel = { Size of Dataset },
			ylabel = {runtime in ms },
			legend style={at={(0.5,-0.4)},
				anchor=south, legend columns=1}]
			
			\addplot[
			color=green!50!red,smooth,tension=0.7,very thick
			] file[skip first] {ABY_use_case2_system_time.dat};
			\addlegendentry{ABY3}

			\addplot[
			color=blue!50!cyan,smooth,tension=0.7,very thick
			] file[skip first] {conclave_use_case2_loca.dat};
			\addlegendentry{Conclave}
		\end{axis}
		
	\end{tikzpicture}
	\caption{Comparison between Conclave and ABY3 implementation in our second use case}
\end{figure}

%\begin{tikzpicture}
%	\tikzset{lines/.style={draw=white},}
%	\pie[color={purple, red, yellow, blue, green},sum=auto, after number=,text=legend,every only number node/.style={text=black},style={lines}]{10/A,20/B,30/C,10/D}
%	\pie[pos={8,0},color={purple, red, yellow, blue, green},sum=auto, after number=,every only number node/.style={text=black},style={lines}]{10/,20/,30/,10/}
%\end{tikzpicture}


\section{Use Case Four}
As we described in Section 6.1 our fourth use case has an identical semantic to our first use case but we make use of Conclave's trust annotations, that allow the leakage of some of the input data, which allows Conclave to apply optimizations that speed up the computation. Therefore we focus on a comparison between the performance of Conclave in the first and fourth use case. At the end out of the section we also compare the first implementation of ABY3 to the implementation of use case four with Conclave to see if Conclave yields comparable result if it has the advantage of its trust annotations.
 A visualisation of our comparison can found in Figures 6.6 and 6.7. 
\paragraph{Comparison between use case One and use case Four }
With the usage of trust annotation we have observed an significant improvement in speed and efficiency. We have been able to evaluate the query for input-sizes of up to 3000 rows, which are 6 times larger then in use case one. For input-sizes larger then 3000 Conclave tends to crash because of internal errors that are outside of our control. Therefore we were unable to obtain result for larger inputs.  For an input of size 500 in use case four, Conclave needs less then 5 minutes which more then 80 times faster then the 8 hours required in use case one. For an input of size 3000, Conclave is able to compute the correct result in less then 25 minutes. From the perspective of memory requirement, the situation is very similar, as Conclave requires less memory in use case four then in use case one, for all input-sizes we observed. 

\paragraph{Comparison between ABY3 and Conclave}
It is notable that even with the unfair advantage of leaking some input data, Conclave is not able to yield similar performance to ABY3. The implementation of use case one with ABY3, which has no such advantage, is still multiple orders of magnitude faster then Conclave in use case four and also requires significantly less memory. On average ABY3 is more then 360 times faster. The difference in speed becomes more significant with larger input-sizes, in the extreme case of input-size 3000 ABY3 is more then 1300 times faster then Conclave. 


\begin{figure}[H]
	\begin{tikzpicture}
		\begin{axis}[
			xmin = 0, xmax = 3000,
			ymin = 0, ymax = 43099000,
			%xtick = {0,4000,8000,16000,32000,64000,128000,512000,1024000,2048000,4096000,8112000,16224000  },
			ymode= log,
			xtick distance=250 ,
			%ytick distance = 20000,
			grid = both,
			minor tick num = 2,
			major grid style = {lightgray},
			minor grid style = {lightgray!25},
			width = \textwidth,
			height = 0.5\textwidth,
			xlabel = { Size of Dataset },
			ylabel = {runtime in ms },
			legend style={at={(0.5,-0.4)},
				anchor=south, legend columns=1}]
			
			\addplot[
			color=green!50!cyan,smooth,tension=0.7,very thick
			] file[skip first] {Conclave_use_case1_loca.dat};
			\addlegendentry{Conclave use case1}
			
			\addplot[
			color=black!50!black,smooth,tension=0.7,very thick
			] file[skip first] {conclave_use_case4_loca.dat};
			\addlegendentry{Conclave use case4}
		\end{axis}
		
	\end{tikzpicture}
\caption{Runtime of Conclave in use case one and use case two}
\end{figure}
\begin{figure}[H]
	\begin{tikzpicture}
		\begin{axis}[
			xmin = 0, xmax = 3000,
			ymin = 0, ymax = 58032576,
			%xtick = {0,4000,8000,16000,32000,64000,128000,512000,1024000,2048000,4096000,8112000,16224000  },
			xtick distance= 250,
			ymode= log,
			%ytick distance = 20000,
			grid = both,
			minor tick num = 2,
			major grid style = {lightgray},
			minor grid style = {lightgray!25},
			width = \textwidth,
			height = 0.5\textwidth,
			xlabel = { Size of Dataset },
			ylabel = {Memory in kb },
			legend style={at={(0.5,-0.4)},
				anchor=south, legend columns=1}]
			
			
			\addplot[
			color=blue!50!cyan,smooth,tension=0.7,very thick
			] file[skip first] {space_conclave_1.dat};
			\addlegendentry{Conclave use case one}
			\addplot[
			color=black!50!black,smooth,tension=0.7,very thick
			] file[skip first] {space_conclave_4.dat};
			\addlegendentry{Conclave use case four}
		\end{axis}
		
	\end{tikzpicture}
	\caption{Comparison between Conlave's space consumption in use case one and use case four}
\end{figure}

\section{Conclusion}
 In our work we have implemented several different use cases for the secure multiparty computation frameworks ABY3, SMCQL, Conclave. We had to overcome problems like missing documentation and complex installation procedures. However, we have able to overcome them and therefore demonstrated the feasibility of the task. 
 In particular, the reasons for problems like the lack of documentation are not inherent properties of multiparty computation, but limited time and resources of the frameworks authors. From the frameworks we used in our tests, only ABY3 was able to demonstrate the capability to handle input sizes of practical relevance, as ABY3 has demonstrated that it is able handle millions of input rows an a moderate hardware. Conclave and SMCQL could only handle a few hundred or thousand input rows. 







%%
\bibliography{references}

\appendix

\end{document}
